% !TEX root = ../main.tex

\section{Related Work}
\label{sec:related_work}
% ==================================================


\subsection{Edge--Cloud VSLAM Transmission}
\label{subsec:edge_cloud_transmission}

Visual SLAM has evolved from feature-based and direct pipelines to learning-based systems with stronger robustness. ORB-SLAM2 and ORB-SLAM3 build sparse maps from repeatable local features and optimize camera poses through bundle adjustment~\cite{murartal2017orbslam2,campos2021orbslam3}. Direct and semi-direct systems such as LSD-SLAM, DSO, and SVO exploit photometric consistency~\cite{engel2014lsdslam,engel2018dso,forster2014svo}, while learning-based systems such as DROID-SLAM introduce dense recurrent optimization across monocular, stereo, and RGB-D inputs~\cite{teed2021droidslam}. These advances improve localization and mapping capability, but they also increase computation and memory demand, motivating edge--cloud VSLAM for resource-constrained robots.

Early collaborative SLAM and cloud robotics systems move expensive mapping, optimization, or reconstruction modules from the robot to external computation resources~\cite{rabaud2011c2tam,kuffner2010cloud,shankar2020fogros}. Recent reviews identify transmission latency and data volume as central bottlenecks, and highlight asynchronous collaboration, submap allocation, and compact representation transmission as important directions~\cite{chen2024reviewcloudedge}. Submap-based edge--cloud VSLAM reduces transmission frequency by offloading selected image segments or submap construction to the cloud~\cite{chen2024cloudlearning,chen2024implicittransmission}. Implicit-representation transmission and VC-SLAM further reduce transferred map data by sending compact neural representations instead of dense point clouds or meshes~\cite{chen2024implicittransmission,huo2025vcslam}, while cross-representation methods study data association and adaptive representation selection for VSLAM~\cite{chen2024crossdata,chen2025xrepslam}.

These works show that edge--cloud VSLAM must reduce communication demand, but they mainly address task allocation, submap transmission, map representation, or data association after visual observations have already been selected. Conventional codecs such as H.264/AVC, H.265/HEVC, and H.266/VVC reduce image streams with rate--distortion objectives designed mainly for human viewing~\cite{sullivan2012h264,sullivan2012hevc,bross2021vvc}. Learned and generative compression exploit neural priors under limited bit budgets~\cite{balle2018variational,minnen2018joint,agustsson2020scale,mentzer2020highfidelity}, and semantic communication emphasizes transmitting task-relevant meaning rather than raw symbols~\cite{qin2021semantic,luo2022semantic}. HVM-Codec addresses the preceding VSLAM-oriented visual transmission problem: encoding the edge image stream itself as sparse image--text visual memory before cloud-side VSLAM, and evaluating that representation by downstream VSLAM trajectory quality.

\subsection{Sparse Visual Memory Coding for Navigation}
\label{subsec:sparse_visual_memory}

Sparse visual memory provides an engineering analogy for replacing dense frame transmission with compact navigation cues. Cognitive studies on event perception suggest that continuous experience can be segmented into events, with boundaries related to changes in prediction, attention, and memory~\cite{zacks2001event,zacks2007event}. Mental-model theory also emphasizes simplified internal representations that preserve structural relationships rather than exhaustive sensory details~\cite{johnsonlaird1983mental}. For VSLAM-oriented visual transmission, this analogy motivates retaining boundary frames, motion transitions, visual novelty anchors, and compact language memory instead of every intermediate frame.

Computer vision has studied related mechanisms through temporal segmentation, change-point detection, and event boundary detection. Generic Event Boundary Detection detects taxonomy-free boundaries according to human-perceived event changes~\cite{shou2021gebd}. Online Bayesian change-point detection provides a statistical foundation for streaming segmentation~\cite{adams2007bayesian}, and streaming representation learning organizes visual inputs into event-level structures~\cite{streamer2023}. Vision-language models such as CLIP and BLIP-2 connect visual content with compact language-level semantics~\cite{radford2021clip,li2023blip2}. In VSLAM, X-RepSLAM uses VLM-derived image-sequence information to guide adaptive representation switching and prompt-driven task customization~\cite{chen2025xrepslam}.

These works provide useful ingredients for segmentation and semantic representation, but they are not designed as a transmissible memory format for edge--cloud VSLAM. Event segmentation typically aims at video understanding or retrieval, and vision-language representations usually describe or select visual content rather than encode a sequence for later task-usable reconstruction. HVM-Codec uses motion memory and visual novelty anchors not only to segment the stream, but also to form image--text visual memory whose boundary frames and language memory are decoded and evaluated by downstream VSLAM.

\subsection{World Models for Visual Sequence Reconstruction}
\label{subsec:world_model_decoding}

World models provide learned priors for predicting or imagining visual continuity from sparse conditions. Ha and Schmidhuber revived the term ``world model'' by learning latent environment dynamics for agents that can simulate possible futures~\cite{ha2018worldmodels}. The Dreamer family demonstrates that latent world models can support decision-making and control through imagined rollouts~\cite{hafner2019planet,hafner2020dreamer,hafner2023dreamerv3}. JEPA-style predictive representation learning emphasizes predicting abstract latent representations rather than raw pixels~\cite{lecun2022path,assran2024vjepa}, and recent surveys categorize world models by their roles in constructing internal representations and predicting future states~\cite{ding2025worldmodelsurvey}.

Large-scale video generation is one visual realization of world models. Image-to-video and text-to-video diffusion models synthesize temporally coherent sequences from sparse visual or language conditions~\cite{blattmann2023svd,ho2022videodiffusion,villegas2022phenaki}. Sora, Genie, Cosmos, and related systems indicate that large video models can capture aspects of spatiotemporal consistency, controllability, and visual imagination, while still having limitations as physical simulators~\cite{brooks2024sora,bruce2024genie,nvidia2025cosmos}. In robotics and embodied AI, world models have been explored for learning from imagined futures, generating robot action videos, and supporting planning~\cite{yang2023unipi,du2023learning,wu2023daydreamer}. Target-Bench is especially relevant because it evaluates generated future videos for mapless path planning and compares recovered camera motion with SLAM-based ground truth trajectories~\cite{wang2025targetbench}.

Generic world models and video generators mainly emphasize future prediction, visual plausibility, realism, diversity, prompt adherence, or planning. VSLAM-oriented reconstruction instead requires stable viewpoints, geometric continuity, and trackable visual textures. Efficient low-rank adaptation such as LoRA makes it possible to bias a pretrained decoder toward a target domain without retraining the full model~\cite{hu2022lora}. HVM-Codec therefore treats the cloud module as an image--text conditioned world-model decoder constrained by sparse visual memory. Its output is not expected to be pixel-identical to the original stream; it is evaluated by task-usable reconstruction and downstream VSLAM trajectory quality.
